\documentclass[11pt, a4paper]{article}

% --- Packages ---
\usepackage[margin=2.5cm]{geometry}
\usepackage{amsmath}
\usepackage{amssymb}
\usepackage{parskip}
\usepackage{hyperref}
\usepackage{microtype}
\usepackage{booktabs}
\usepackage{enumitem}

% --- Hyperref setup ---
\hypersetup{
    colorlinks=true,
    linkcolor=black,
    urlcolor=blue,
    citecolor=black
}

% --- Title ---
\title{\textbf{Quantitative Dynamic Macroeconomics}\\[0.4em]
\large Paper Summary: Clarida, Gal\'{i}, and Gertler (1999)}
\author{Konstantin Gantert}
\date{\today}

% ---------------------------------------------------------------
\begin{document}
\maketitle

\noindent Clarida, R., Gal\'{i}, J., and Gertler, M.\ (1999). The Science of
Monetary Policy: A New Keynesian Perspective. \textit{Journal of Economic
Literature}, 37(4), 1661--1707.

\tableofcontents
\bigskip

% ---------------------------------------------------------------
\section{Introduction}

By the late 1980s there was broad agreement that monetary policy has important
consequences for aggregate activity, and the theoretical frameworks available
for policy analysis had improved considerably. CGG position their paper as a
synthesis of what this progress produced.

The central question is normative: what should a central bank do, and why?
The goal is to derive principles that are \textbf{robust} --- holding across
a wide range of model specifications within the New Keynesian class rather
than depending on particular parameter choices. Results that are highly
model-specific are of limited use for practical policy guidance. The paper's
ambition is therefore to identify principles that survive across the NK class
and to assess how well they square with actual central bank behaviour.

The paper deliberately restricts its scope to models where nominal price
rigidities are the key friction generating monetary non-neutrality. The
principles CGG derive are valid within this class of models, and the project
model belongs to exactly this class.

% ---------------------------------------------------------------
\section{A Baseline Framework}

The model consists of two equations familiar from lecture.

\textbf{The IS curve:}
\begin{equation}
    x_t = -\varphi[i_t - E_t\pi_{t+1}] + E_t x_{t+1} + g_t
\end{equation}

\textbf{The New Keynesian Phillips curve:}
\begin{equation}
    \pi_t = \lambda x_t + \beta E_t\pi_{t+1} + u_t
\end{equation}

The general equilibrium rests on Walrasian labor markets and monopolistically
competitive goods markets. A key difference from the traditional IS/LM
framework is that all behavioral relationships are derived from optimising
households and firms. As a result, \textbf{beliefs about future policy
directly affect current inflation and the output gap} through the
forward-looking structure of both equations. Expected future interest rates
matter as much as current interest rates.

Both equations are written in terms of \textbf{gap variables}. The output
gap $x_t = y_t - z_t$ measures actual output relative to the flexible-price
benchmark $z_t$. Inflation is driven by this gap, not by the level of output.
A boom driven by a genuine rise in potential output requires no policy
response. A boom that opens a positive output gap does.

The \textbf{policy objective} formalises this:
\begin{equation}
    \mathcal{L} = E_0 \sum_{t=0}^{\infty} \beta^t
    \left[\alpha x_t^2 + \pi_t^2\right]
\end{equation}

The central bank minimises squared deviations of the output gap and inflation
from target. The relative weight $\alpha$ on output gap variance is a
structural parameter that depends on the degree of price stickiness: stickier
prices imply a larger weight on output stabilisation. The monetary policy
instrument is the nominal interest rate --- no LM curve is needed since money
markets clear implicitly.

% ---------------------------------------------------------------
\section{Optimal Monetary Policy without Commitment}\label{sec:discretion}

This section derives optimal policy under \textbf{discretion} --- the
empirically relevant case since no major central bank makes binding
commitments. Discretion is defined by two conditions: (a) future inflation
and the output gap are not affected by today's central bank actions; and (b)
the central bank cannot directly manipulate private sector expectations.

\subsection*{The Output-Inflation Trade-Off}

The central bank's optimality condition under discretion is:
\begin{equation}
    x_t = -\frac{\lambda}{\alpha}\pi_t
\end{equation}

This ``lean against the wind'' condition says the central bank should
contract demand whenever inflation is above target. A trade-off between
output and inflation variance exists \textbf{if and only if cost-push shocks
are present}. Without cost-push shocks ($u_t = 0$), the central bank can
simultaneously stabilise both inflation and the output gap --- the divine
coincidence. With cost-push shocks, reducing inflation requires contracting
demand below potential, making the trade-off unavoidable.

\subsection*{Inflation Targeting is Generally Gradual}

Returning inflation immediately to target --- extreme inflation targeting ---
is only optimal in two limiting cases: when there are no cost-push shocks
($\sigma_u = 0$), or when the central bank places zero weight on output
($\alpha = 0$). In the general case with both cost-push shocks and
$\alpha > 0$, gradual convergence of inflation to target is optimal. The
efficient policy frontier in output-inflation volatility space is downward
sloping and convex: moving toward lower inflation variance always comes at
the cost of higher output gap variance when cost-push shocks are present.

\subsection*{The Taylor Principle}

The optimal interest rate rule takes the form:
\begin{equation}
    i_t = \gamma_\pi E_t\pi_{t+1} + \frac{1}{\varphi}\,g_t,
    \qquad \gamma_\pi > 1
\end{equation}

where $g_t$ is the demand shock from the IS curve --- a function of expected
changes in government purchases relative to expected changes in potential
output. Since the optimal policy perfectly offsets demand shocks, $g_t$
enters the interest rate rule directly. The key result is $\gamma_\pi > 1$:
the nominal rate must respond more than one-for-one with expected inflation
so that the \textit{real} interest rate rises when inflation rises. Violating
this --- setting $\gamma_\pi \leq 1$ --- permits self-fulfilling inflationary
spirals. CGG estimate this coefficient for the US Federal Reserve and find it
was below one in the pre-Volcker era (1960--1979) and above one in the
Volcker--Greenspan era (1979--1996), which is revisited in detail in
Section~\ref{sec:simple_rules}.

\subsection*{Demand versus Supply Shocks}

The optimal policy response depends critically on the source of the shock:
\begin{itemize}[noitemsep]
    \item \textbf{Demand shocks:} fully offset. Raising rates closes the
          output gap and simultaneously stabilises inflation --- no trade-off.
    \item \textbf{Supply shocks to potential output:} fully accommodate.
          A rise in potential output raises demand in an exactly offsetting
          way, leaving the output gap unchanged. No rate response is needed.
    \item \textbf{Cost-push shocks:} balance. These cannot be fully offset
          without causing an output contraction. This is where the trade-off
          bites.
\end{itemize}

% ---------------------------------------------------------------
\section{Credibility and the Gains from Commitment}

\subsection*{The Classic Inflation Bias}

When the central bank targets output above potential ($k > 0$), discretion
leads to permanently higher inflation with no output gain:
\begin{equation}
    \pi_t^k = \pi_t + \frac{\alpha}{\lambda}\,k
\end{equation}

Rational agents anticipate the central bank's temptation to expand demand
and price accordingly, eliminating any output benefit. The bias is larger
when the central bank cares more about output ($\alpha$ high) and when the
Phillips curve is flatter ($\lambda$ small). Rogoff's (1985) solution ---
appointing a central banker with a stronger distaste for inflation,
effectively lowering $\alpha$ --- reduces the bias.

\subsection*{Gains from Commitment even with $k = 0$}

The more novel result is that commitment improves outcomes even when the
central bank has no output ambitions beyond the natural level. Because
current inflation depends on expected future output gaps through the
forward-looking Phillips curve, a central bank that credibly commits to
fighting inflation in the future immediately reduces current inflation without
requiring a large current recession. The welfare cost of reducing inflation
falls --- disinflating requires a smaller output contraction under commitment
than under discretion. Formally, the optimal committed rule implies an
effective weight on inflation of (where $\rho$ is the persistence of the cost-push shock):
\begin{equation}
    \alpha^c \equiv \alpha(1-\beta\rho) < \alpha
\end{equation}

The central bank behaves as if it cares more about inflation than it actually
does. This outcome can be replicated under discretion by appointing a
sufficiently conservative central banker.

\subsection*{The Unconstrained Optimum}

The globally optimal committed policy satisfies a \textit{difference rule}
rather than a level rule:
\begin{equation}
    x_{t+i} - x_{t+i-1} = -\frac{\lambda}{\alpha}\pi_{t+i},
    \qquad i \geq 1
\end{equation}

The central bank adjusts the \textit{change} in the output gap in response
to inflation, continuing to contract demand even after inflation has begun
falling. This history dependence is the source of commitment's value but
makes the rule harder to implement in practice.

% ---------------------------------------------------------------
\section{Practical Complications}

\subsection*{Imperfect Information}

When the central bank cannot observe the current state of the economy in real
time --- because data are delayed or key variables like potential output are
unobservable --- optimal policy rules remain valid in their
certainty-equivalent form but must be expressed in terms of forecasts:
\begin{equation}
    E\{x_t \,|\, \Omega_t\} = -\frac{\lambda}{\alpha}\,E\{\pi_t \,|\, \Omega_t\}
\end{equation}

Since monetary policy affects inflation with a significant lag --- empirical
estimates suggest output responds within six to nine months and inflation
within around eighteen months, though these vary across models and contexts
--- the central bank should target its inflation forecast rather than current
inflation. This is the practical case for inflation forecast targeting.

\subsection*{Interest Rate Smoothing}

Central banks adjust rates far more gradually than optimal rules predict. A
fitted Taylor rule for the US yields smoothing parameters $\rho_i$ around
0.8--0.9. Two explanations are offered: parameter uncertainty in the Brainard
(1967) tradition induces caution, and concerns about disrupting financial
markets may prevent sharp interest rate movements. Neither explanation is
fully satisfying and CGG flag this as an important unresolved issue.

% ---------------------------------------------------------------
\section{Endogenous Persistence: The Hybrid Model}

Section~6 extends the baseline to allow for endogenous persistence in output
and inflation, motivated empirically. The key new result is that with
inflation inertia, monetary policy affects not just how far inflation deviates
from target but also \textbf{how quickly it returns}. The efficient policy
frontier shifts unfavourably: any given degree of inflation persistence
worsens the output-inflation trade-off.

The important message is the \textbf{robustness} finding: all four main
results from Section~\ref{sec:discretion} --- the trade-off, gradual
targeting, the Taylor principle, and the differential response to demand
versus supply shocks --- survive the extension to the hybrid model. The
qualitative prescriptions are unchanged even when the model is substantially
generalised. This is the strongest support CGG offer for their claim that
there is a science of monetary policy.

% ---------------------------------------------------------------
\section{Simple Rules for Monetary Policy}
\label{sec:simple_rules}

\subsection*{The Taylor Rule}

Taylor (1993) proposed:
\begin{equation}
    i_t^* = \bar{r} + \pi^* + \gamma_\pi(\pi_t - \pi^*) + \gamma_x x_t
\end{equation}

with $\gamma_\pi = 1.5$ and $\gamma_x = 0.5$. CGG assess this rule against
the normative principles established in earlier sections and find it passes
on all counts: it satisfies the Taylor principle, calls for gradual
adjustment of inflation, responds to the output gap rather than the level of
output, and implicitly accommodates supply shocks while offsetting demand
shocks.

CGG then estimate a forward-looking variant with interest rate smoothing for
two subsamples of US monetary policy history:

\begin{center}
\begin{tabular}{lccc}
\toprule
& $\gamma_\pi$ & $\gamma_x$ & $\rho_i$ \\
\midrule
Pre-Volcker (1960--1979)       & 0.83 & 0.27 & 0.68 \\
                               & (0.07) & (0.08) & (0.05) \\[0.3em]
Volcker--Greenspan (1979--1996) & 2.15 & 0.93 & 0.79 \\
                               & (0.40) & (0.42) & (0.04) \\
\bottomrule
\end{tabular}
\end{center}

The pre-Volcker estimate violates the Taylor principle ($\gamma_\pi < 1$),
meaning the Fed was accommodating rises in expected inflation rather than
raising real rates to fight them. CGG argue this directly explains the
persistent high inflation of the 1970s. The post-1979 estimate is strongly
above one, consistent with the implicit inflation targeting that characterises
good policy. The Great Moderation --- the sharp reduction in macroeconomic
volatility from the mid-1980s onward --- is interpreted as the consequence of
this regime shift. While alternative explanations exist --- including good
luck and structural changes in the economy --- the policy regime shift under
Volcker is widely cited as a leading factor.

\subsection*{Indeterminacy}

When $\gamma_\pi \leq 1$, a rise in inflation expectations lowers the real
rate, stimulates demand, and raises inflation further --- a self-fulfilling
spiral. The pre-Volcker rule permitted exactly this. Satisfying the Taylor
principle is therefore not merely a normative prescription but a condition for
ruling out destabilising sunspot dynamics.

\subsection*{Further to Consider}

CGG also assess two alternative target rules.

\textbf{Inflation targeting} --- committing to a specific inflation target
and adjusting policy gradually --- is broadly consistent with the optimal
policy from earlier sections and provides a clear nominal anchor that
disciplines discretion.

\textbf{Price level targeting} is found unattractive: any overshoot of the
price level must be reversed, requiring future deflation to correct past
inflation mistakes, which introduces unnecessary output volatility.
Moreover, measurement error in the price level can trigger contractionary
policy even when none is warranted.

\textbf{Nominal GDP targeting} applies equal weight to inflation and real
output growth, which is generally not optimal. A supply shock raising
potential output triggers a contractionary response even though none is
needed, and permanent real GDP growth means there is no precise nominal
anchor.

% ---------------------------------------------------------------
\section{Research Agenda}

CGG close with six open questions: the sources of inflation persistence, open
economy extensions, the zero lower bound on nominal interest rates, interest
rate smoothing, the role of financial stability in monetary policy, and
private sector learning during regime transitions.

In the twenty-five years since publication, each of these has generated an
enormous literature. The zero lower bound became the central challenge of
monetary policy after 2008. Financial stability moved to the core of central
bank mandates after the Global Financial Crisis. Inflation persistence
returned as a live debate during the post-COVID inflation surge of
2021--2023. Students interested in monetary economics will encounter this
literature in a research master or PhD programme. For the purposes of this
course, the research agenda is a reminder that the CGG framework is a
deliberately simplified lens --- powerful, but not the final word.

\end{document}
